\chapter{97}

\section{Mittwoch, 12.
September}



Manuela und Ava zogen noch vor Sonnenaufgang los. Vollgepackt mit all
dem Notwendigen.

«Wir stinken bald wie offene Müllsäcke, David und ich, und sind knapp am
Verhungern,» hatte Mazi am Vorabend hinunter gemeldet. Er hatte sich
gekonnt verzweifelt gegeben, aber darauf übermütig ins Handy gelacht.

Trotz Mazis humorvoller Einlage hatte sich Joh schäbig gefühlt.

«Wie kann es sein, B.B., dass ich mich so krass in Louis’ Drama
verliere? Manchmal fürchte ich, die beiden oben zu vergessen. Shit.»

Ben erhob sich vom Sofa. Mit Notizbuch und Stift näherte er sich Joh.
Der sass bereits am grossen Tisch und rückte eine Kaffeetasse auf seine
gegenüberliegende Seite. Ben setzte sich. Er fuhr sich mit den Fingern
durchs Haar und seufzte.

«Hm, ich such gerade nach Trost. Ich denk mir, grosse Katastrophen
generieren grosse Folgen.»

Dann richtete er sich gerade auf.

«Aber eigentlich ist es ein gutes Zeichen. Du vertraust darauf, dass
sie’s oben richten. Und das tun sie doch ganz gut, hm?»

«Auch wieder wahr, B.B.,» und Johs Stimme gewann an Bestimmtheit, «dann
also zu unserem Vorgehen. Sprechen wir unseren Plan nochmal durch, bevor
Louis da ist, okay?»

\sectionbreak

Auch wenn Joh seinen Freund gut kannte, hatte er sich vergewissern
müssen, ob er Ben richtig verstanden hatte.
Sie hatten am Vortag lange mit Louis zusammen gesessen. 
«Ich kann mir deine Unterstützung nicht leisten, Ben,» hatte Louis 
gesagt und dabei ausgesehen, wie ein geprügelter Hund.
Bens Antwort hatte auf sich warten lassen. Sein Zögern hatte in
Windeseile auch Fragen in Joh ausgelöst.
Was, wenn Bens Unterstützung zum juristischen Grossaufgebot wurde? Wer
wusste schon, was da noch alles ans Licht kam? Was, wenn Ben Louis’
«Fall» als offizielles Mandat behandeln musste und das Ganze uferlos
lange dauern würde?
Über sich häufende Gerichtskosten hatte er kaum noch nachdenken mögen.
Noch viel weniger über ein eventuelles Anwaltshonorar für Ben.

Louis’ Lügengeschichte setzte Joh mehr zu, als ihm lieb war. Darüber unterhielt er sich
später alleine mit Ben. 

Als Ben sich verabschiedet hatte, blieb Joh alleine auf dem Sofa sitzend zurück. Er liess seinen Blick durch den Wohnraum gleiten. Dann hielt er inne und starrte die kleine Vitrine an. Sie hing an der Seitenwand gegenüber und schien ihn listig anzugrinsen. Hinter der schmalen Glaswand sah er sie. Eine einzelne, elegant ausgestellte Zigarette, die er über viele Jahre kaum mehr beachtet hatte.  
Ein Geschenk von Manuela, als er ihr seinen Rauchstopp angekündigt hatte. Sie stehe als Symbol für das Ende seiner Sucht, hatte Manuela damals gesagt, ihn fest an sich gedrückt «für deine Sehnsucht.» Dann hatte sie gelacht und angefügt, «beende dein Sehnen, leb das Leben.»
Er musste sich zurückhalten. Am liebsten hätte er das abgeriegelte Glastürchen eingeschlagen. Hätte er sich die Zigarette geschnappt, angezündet und nach einem tiefen Atemzug auf Entspannung gewartet. Manuela kam langsam und mit einem Tablett Teetassen aus der Küche. Sie setzte sich wortlos neben Joh. Er klopfte sich ärgerlich auf die Schenkel.

«Hätte ich die Formalitäten rund um Louis` Anstellung sofort erledigt, wäre er augenblicklich aufgeflogen.»

Manuela schaute ihn nachdenklich an.

«Du bist zu streng mit dir, Joh. Wann, bitte, hättest du denn Zeit dazu gehabt? Die letzten Wochen waren turbulent,» sie atmete tief ein, «nimm's als Schicksal - ich glaube, wir lernen gerade etwas Wichtiges - ,» sie küsste ihn auf die Wange und schmunzelte, «und ein Mister Perfect an meiner Seite wäre mir überdies zu langweilig.»

Schliesslich hatte sich Joh beruhigt, lag aber dennoch stundenlang wach. Er versuchte nach wie vor zu verstehen, warum ihn selbst rotzfreche Lügen seiner «Schützlinge» nicht entfernt so tief
erschütterten, wie Louis’ Identitätsschwindel. 
Immerhin konnte er auf zwei seiner ehemaligen Jungs bauen. Der Gedanke, dass sie ihn unterstützen wollten, hellte ihn auf. Ohne Nachzudenken hatten sie zugesagt, die bevorstehende  \emph{Schafscheid}  zu begleiten. Vielleicht, hatte Manuela gemeint, sei es doch ganz gut, wenn Louis während der Tage die \emph{Sunnewaid} in Schuss halte. 








